\textbf{Phần 1. Trắc nghiệm lựa chọn}
\setcounterpageref{ex}{0}

\begin{ex} %[3D2N14]
Một hình tròn được chia thành 4 phần bằng nhau, đã tô màu 1 phần. Phân số chỉ phần đã tô màu là
\choice
{Một phần ba}
{\True Một phần tư}
{Một phần năm}
{Một phần hai}
\loigiai{
Đã tô màu $\dfrac{1}{4}$ (một phần tư) hình tròn.\\
\textbf{Đáp án của câu hỏi là B}
}
\end{ex}

\begin{ex} %[3D2N14]
$\dfrac{1}{2}$ còn được gọi là
\choice
{\True Một nửa}
{Một phần tư}
{Một phần ba}
{Gấp đôi}
\loigiai{
$\dfrac{1}{2}$ gọi là một nửa.\\
\textbf{Đáp án của câu hỏi là A}
}
\end{ex}

\begin{ex} %[3D2H14]
$\dfrac{1}{3}$ của $18\text{ kg}$ là
\choice
{$5\text{ kg}$}
{$7\text{ kg}$}
{\True $6\text{ kg}$}
{$8\text{ kg}$}
\loigiai{
$18 : 3 = 6\text{ kg}$.\\
\textbf{Đáp án của câu hỏi là C}
}
\end{ex}

\begin{ex} %[3D2H14]
$\dfrac{1}{5}$ của $35\text{ l}$ là
\choice
{$6\text{ l}$}
{$8\text{ l}$}
{$5\text{ l}$}
{\True $7\text{ l}$}
\loigiai{
$35 : 5 = 7\text{ l}$.\\
\textbf{Đáp án của câu hỏi là D}
}
\end{ex}

\begin{ex} %[3D2H14]
$\dfrac{1}{6}$ của $42\text{ m}$ là
\choice
{$6\text{ m}$}
{\True $7\text{ m}$}
{$8\text{ m}$}
{$9\text{ m}$}
\loigiai{
$42 : 6 = 7\text{ m}$.\\
\textbf{Đáp án của câu hỏi là B}
}
\end{ex}

\begin{ex} %[3D2H14]
$\dfrac{1}{8}$ của $64\text{ quả cam}$ là
\choice
{$7$ quả}
{$9$ quả}
{\True $8$ quả}
{$6$ quả}
\loigiai{
$64 : 8 = 8$ quả cam.\\
\textbf{Đáp án của câu hỏi là C}
}
\end{ex}

\begin{ex} %[3D2V14]
Một lớp có 36 học sinh, trong đó $\dfrac{1}{4}$ số học sinh là học sinh giỏi. Số học sinh giỏi của lớp đó là
\choice
{\True $9$ học sinh}
{$8$ học sinh}
{$10$ học sinh}
{$12$ học sinh}
\loigiai{
$36 : 4 = 9$ học sinh.\\
\textbf{Đáp án của câu hỏi là A}
}
\end{ex}

\begin{ex} %[3D2V14]
Mẹ có 45 quả táo, mẹ đã biếu bà $\dfrac{1}{5}$ số táo đó. Số quả táo mẹ đã biếu bà là
\choice
{$8$ quả}
{$10$ quả}
{$7$ quả}
{\True $9$ quả}
\loigiai{
$45 : 5 = 9$ quả táo.\\
\textbf{Đáp án của câu hỏi là D}
}
\end{ex}

\begin{ex} %[3D2V14]
Một bao gạo nặng $54\text{ kg}$. Người ta đã bán đi $\dfrac{1}{6}$ bao gạo. Số gạo còn lại trong bao là
\choice
{$43\text{ kg}$}
{\True $45\text{ kg}$}
{$47\text{ kg}$}
{$40\text{ kg}$}
\loigiai{
Đã bán: $54 : 6 = 9\text{ kg}$.\\
Còn lại: $54 - 9 = 45\text{ kg}$.\\
\textbf{Đáp án của câu hỏi là B}
}
\end{ex}

\begin{ex} %[3D2V14]
Một đoạn dây dài $72\text{ cm}$. Lan cắt lấy $\dfrac{1}{9}$ đoạn dây để làm nơ. Đoạn dây Lan đã cắt dài là
\choice
{\True $8\text{ cm}$}
{$7\text{ cm}$}
{$9\text{ cm}$}
{$6\text{ cm}$}
\loigiai{
$72 : 9 = 8\text{ cm}$.\\
\textbf{Đáp án của câu hỏi là A}
}
\end{ex}

\begin{ex} %[3D2V14]
Có 48 bông hoa, cắm $\dfrac{1}{8}$ số hoa vào lọ thứ nhất. Lọ thứ nhất có số bông hoa là
\choice
{$5$ bông}
{$7$ bông}
{$8$ bông}
{\True $6$ bông}
\loigiai{
$48 : 8 = 6$ bông hoa.\\
\textbf{Đáp án của câu hỏi là D}
}
\end{ex}

\begin{ex} %[3D2C14]
Một đàn gà có 40 con. Biết $\dfrac{1}{5}$ số gà là gà trống, còn lại là gà mái. Số gà mái có trong đàn là
\choice
{$30$ con}
{$34$ con}
{\True $32$ con}
{$28$ con}
\loigiai{
Gà trống: $40 : 5 = 8$ con.\\
Gà mái: $40 - 8 = 32$ con.\\
\textbf{Đáp án của câu hỏi là C}
}
\end{ex}

\textbf{Phần 2. Trắc nghiệm đúng sai}
\setcounterpageref{ex}{0}

\begin{ex} %[3D2H14]
Xét tính đúng sai về phần bằng nhau của một số:
\choiceTF
{\True $\dfrac{1}{2}$ của $20$ là $10$.}
{\True $\dfrac{1}{3}$ của $27$ là $9$.}
{$\dfrac{1}{4}$ của $32$ là $7$.}
{\True $\dfrac{1}{4}$ của $32$ là $8$.}
\loigiai{
\begin{itemize}
    \item $20 : 2 = 10$: Đúng.
    \item $27 : 3 = 9$: Đúng.
    \item $32 : 4 = 8 \neq 7$. Ý 3 Sai, Ý 4 Đúng.
\end{itemize}
\begin{itemize}
    \item a) \textbf{Đúng}.
    \item b) \textbf{Đúng}.
    \item c) \textbf{Sai}.
    \item d) \textbf{Đúng}.
\end{itemize}
}
\end{ex}

\begin{ex} %[3D2V14]
Một cuộn vải dài $63\text{ m}$. Cửa hàng đã bán được $\dfrac{1}{7}$ cuộn vải. Xét tính đúng sai:
\choiceTF
{\True Số vải đã bán là $9\text{ m}$.}
{\True Số vải còn lại là $54\text{ m}$.}
{Số vải còn lại gấp 5 lần số vải đã bán.}
{\True Số vải còn lại gấp 6 lần số vải đã bán ($54 : 9 = 6$).}
\loigiai{
\begin{itemize}
    \item Đã bán: $63 : 7 = 9\text{ m}$. Đúng.
    \item Còn lại: $63 - 9 = 54\text{ m}$. Đúng.
    \item $54 : 9 = 6$ lần $\neq 5$ lần. Ý 3 Sai, Ý 4 Đúng.
\end{itemize}
\begin{itemize}
    \item a) \textbf{Đúng}.
    \item b) \textbf{Đúng}.
    \item c) \textbf{Sai}.
    \item d) \textbf{Đúng}.
\end{itemize}
}
\end{ex}

\textbf{Phần 3. Trả lời ngắn}
\setcounterpageref{ex}{0}

\begin{ex} %[3D2H14]
Tìm $\dfrac{1}{5}$ của $40\text{ kg}$.
\shortans{8}
\loigiai{
\textbf{Đáp số:} $8$\\
$40 : 5 = 8\text{ kg}$.
}
\end{ex}

\begin{ex} %[3D2H14]
Tìm $\dfrac{1}{7}$ của $49\text{ m}$.
\shortans{7}
\loigiai{
\textbf{Đáp số:} $7$\\
$49 : 7 = 7\text{ m}$.
}
\end{ex}

\begin{ex} %[3D2V14]
Tìm $\dfrac{1}{8}$ của $56\text{ l}$.
\shortans{7}
\loigiai{
\textbf{Đáp số:} $7$\\
$56 : 8 = 7\text{ l}$.
}
\end{ex}

\begin{ex} %[3D2V14]
Tìm $\dfrac{1}{9}$ của $81$ quả táo.
\shortans{9}
\loigiai{
\textbf{Đáp số:} $9$\\
$81 : 9 = 9$ quả táo.
}
\end{ex}

\begin{ex} %[3D2V14]
Một tấm vải dài $36\text{ m}$, người ta đã may áo hết $\dfrac{1}{4}$ tấm vải đó. Hỏi đã may áo hết bao nhiêu mét vải?
\shortans{9}
\loigiai{
\textbf{Đáp số:} $9$\\
$36 : 4 = 9\text{ m}$.
}
\end{ex}

\begin{ex} %[3D2C14]
Một rổ có 30 quả cam, An ăn $\dfrac{1}{6}$ số cam đó, Bình ăn $\dfrac{1}{5}$ số cam ban đầu. Hỏi cả hai bạn đã ăn bao nhiêu quả cam?
\shortans{11}
\loigiai{
\textbf{Đáp số:} $11$\\
An ăn: $30 : 6 = 5$ quả. Bình ăn: $30 : 5 = 6$ quả. Cả hai bạn ăn: $5 + 6 = 11$ quả cam.
}
\end{ex}
